\documentclass[11pt,a4paper]{article}
\usepackage[
  top=2.5cm,
  bottom=2.5cm,
  left=2.5cm,
  right=1.5cm,
]{geometry}
\usepackage[slovak]{babel}
\usepackage[T1]{fontenc}

\usepackage{tabularx}

\def\variant{vv-BVS}
\def\repo{https://github.com/TekMike365/ifj26c.git}
\def\extensions{}

\begin{document}
  \begin{titlepage}
    \vspace*{40mm}

    \centering \Huge
    \textbf{Dokumentácia \texttt{ifj26c}}
    \vspace{1.25em}

    \centering \Large
    \begin{tabularx}{\textwidth}{||r X}
           Varianta: &\variant \\
      Git repozitár: &\repo \\
         Rozšírenia: &\extensions \\
    \end{tabularx}

    \vfill

    \raggedleft \Large
    \begin{tabular}{l l r|r}
      \hline\hline
                       &           &       & body [\%] \\
      Sklenka Michael, &xsklenm00, &281936 & 25 \\
      Makuka Dominik,  &xmakukd00, &281937 & 25 \\

      \bfseries Dedič Ondrej, &\bfseries xdedico00, &\bfseries 282358
      & \bfseries 25 \\

      Blšták Richard,  &xblstar00, &282097 & 25 \\
    \end{tabular}
  \end{titlepage}

  \tableofcontents
  \newpage

  \section{Rozdelenie prác}
  % stručný popis rozdělení práce mezi členy týmu (uveďte výčtem nebo několika
  % větami, kdo a jak se podílel na jednotlivých částech projektu; povinně
  % zdůvodněte odchylky od rovnoměrného rozdělení bodů);

  \section{Spolupráca a kvalita kódu}
  % stručný popis způsobu spolupráce, kultury code review a schvalování PR (max.
  % jednotky vět či odrážek); stručný popis využitých forem testování (max.
  % jednotky vět či odrážek); zvolený styl formátování kódu;

  \section{Architektúra}
  % blokové schéma architektury vašeho řešení, které bude obsahovat jednotlivé
  % moduly překla- dače (např. lexikální, syntaktická, sémantická analýza atp.)
  % a musí zobrazovat i datové toky mezi jednotlivými komponentami –
  % specifikujte, v jakých strukturách (např. tokeny, AST, me- zikód) si
  % komponenty předávají informace. Z diagramu a jeho textového popisu musí být
  % jasné vstupy, výstupy, význam komponent a jejich propojení;

  \section{Lexikálny analyzátor}
  % diagram automatu, který specifikuje lexikální analyzátor (podle konvencí
  % přednášených v IFJ; s případným slovním popisem odchylek od standardního
  % modelu konečného automatu);

  \section{Syntaktický analyzátor}
  % LL gramatiku a z ní vycházející LL tabulku; a precedenční tabulku, podle
  % kterých jste imple- mentovali syntaktický analyzátor;

  \section{Moduly}
  % stručný popis každého modulu (informačně obhacující věta na dvojici .c/.h);

  % TODO: bibliography
  % bibliografické citace všech převzatých zdrojů (vč. kódu, algoritmů);

  \section{Využitie AI}
  % popis využití prostředků AI (viz sekci 4 v zadaní).

\end{document}
